% !TeX root = ../papers/nc-p3.tex
\section{被約四頂点切断の左右核座標の反対対称性}
\label{sec:finite-tetrahedral-kernel-symmetry}
\begin{theorem}\label{thm:finite-tetrahedral-kernel-symmetry}
四辺が有限で全中央支持4の有限Frobenius窓に共通印付き切断 $E$ があり、
その階数一切断 $\Sigma$ が被約四点であるとする。印の接次元が高々1なら、$E$ の
行・列に当たる二位置（以下、奇位置）の間に射影線型な二移動があり、その側の反対対称性が成り立つ。
他の局所因子の接次元・被約性を仮定しない。
\end{theorem}
\begin{proof}[外積の支持と三次式への持上げ]
位置1,3の印からの三射影座標をそれぞれ $l,r$ とする。
四頂点表示を $A=\sum_i a_i u_i u_i^t$, $B=\sum_i b_i u_i u_i^t$ と取り、
$\sum u_i=0$ となるよう四方向を規格化する（係数は二つの偶座標に吸収する）。
どの三方向も独立である。$l_i=u_i^tl$ とすると $Al=0$ より
$a_i l_i$ は四つとも等しい。$Al=Ar=Bl=Br=0$ は点環 $D$ 上の等式である。
奇側が非対称と仮定する。$k=l\times r$ の三成分は定理\ref{thm:finite-opposite-product-gap}より独立である。

$A$ の各列 $v$ は $l,r$ の両方を消すから $v\times k=0$ である。
各局所因子で $A$ の単元成分を持つ列を取り、その列に比例して $k=\eta v$ と書ける。
他の全列に同じ恒等式を使うと $A$ の全二次小行列式は $\eta$ に消される。
従って $k$ の成分の生成イデアルは局所主で、その零化イデアルが定める $Z$ は
$A$ により $\Sigma$ に入る。偶座標から $\mathbb P(E)$ への写像が線型同型であるため、
$Z$ は相異なる $n\le4$ 個の被約点である。$B$ にも同じ議論を適用すると、
$Z$ 上で $A,B$ の列は同じ $v$ の方向にあり、二つの偶座標は同じ頂点に対応する。
$Z$ の各点 $j$ での $k$ は零化イデアルが極大イデアルのベクトル、すなわち
$u_j$ と非零socle元の積である。三成分の独立性から $n=3,4$ である。
印での $A$ は階数 $3-\dim T_pX\ge2$ なので $Z$ は印を含まない。

位置1,3の四座標を $U,W$ とし、単元 $h$ を $h^{-1}W^2=U^2$ と取る。
$c_j$ を上のsocle元の $h^{-1}$ 倍を表す位置1の三次式とする。三次制限が同型なので
これは一意である。局所自明化に従い、$b\in W$ の値を $b(j)$ と書く。このとき
\[
                    \sum_{j\in Z} b(j)\,l_j c_j=0
                    \qquad (b\in W)
\]
が多項式恒等式である。実際 $h^{-1}b k$ の三次持上げは一方では
$\sum_j b(j)c_j u_j$、他方では $l$ と、三つの $h^{-1}b r_\nu$ の二次持上げの列との外積である。
後者と $l$ の内積は多項式として零である。
$W$ の $Z$ 上の値の階数は2以上で、$n$ より小さい。全射なら各非零積 $l_j c_j$ が零になるからである。

階数が $n-1$ なら定数な一つの線型従属の係数は全て非零であり、$l_j c_j$ は
いずれも同じ多項式の非零定数倍である。異なる $l_j$ は比例しないので
\[
                    c_j=\gamma_j H\prod_{i\in Z,\ i\ne j}l_i ,
      \qquad \gamma_j\in k^\times,\quad\deg H=4-n .
\]
\end{proof}

\begin{proof}[長さ三と四の排除]
$n=3$ なら $Z$ の添字を $1,2,3$ とし $H$ は非零一次式である。
隣接座標積 $H a_4$ の支持長は高々5であることを示す。
印以外では四つの $l_i$ のうち少なくとも二つが単元である（三方向ずつ独立である）。
三つの $c_j$ を法とすると、$1,2,3$ のうち二つが単元の場合には $H=0$ となる。
残りの場合は $l_4$ と一つが単元で、他の二つの $l_i$ に $H$ を掛けると零になる。
$a_4 l_4=a_i l_i$ より $H a_4=0$ である。従って印以外で $H a_4$ は三つのsocleの包に入る。
印の環を $k[t]/(t^m)$ とし $m>2$ の場合だけを考える。$l/t$ の剰余は非零なので
$l_i/t$ の少なくとも二つが単元である。$c_j$ はここで零であり、$t^2$ で割った式と
$a_i l_i/t$ の一致を法 $t^{m-2}$ で用いると同じ議論で $Ha_4=0\bmod t^{m-2}$。
これで支持長の上界を得る。補題\ref{lem:finite-adjacent-word-support}から $Ha_4=0$ となるが、
これは有限Frobenius窓に純関係がないことに反する。

$n=4$ で値の階数が3なら $H$ は定数である。印以外で二つの $l_i,l_j$ を単元に取ると、
他の二つの $l$ 座標は $c_j$ 全体のイデアルに入る。
従って $a_i,a_j$ もそのイデアルに入る。これは $Z$ 上だけのsocleイデアルである。
$Z$ の各点では $a$ の四座標のうちただ一つだけが単元であるため、相異なる三つの
$a$ 座標の積は印以外で零である。実際少なくとも一因子が上のsocleに入り、
さらに一因子が極大イデアルに入る。印で $m>3$ なら $t^3$ で割った三積の消滅と
$a_i l_i/t$ の一致から、少なくとも二つの $a_i$ は $t^{m-3}$ の倍となる。
従って相異なる三座標の四つの三積は、印の長さ高々3のイデアルに全て入る。
これは偶位置の三次制限の単射性に反する。

残りは $n=4$ で位置3での四点が共線の場合である。有限窓において、ある位置で
四点が共線で、その前点と後点がそれぞれ射影基底をなすことは不可能である。
ここで射影基底とは四次元線型空間の基底の四点を指す。
実際、前・後の点をそれぞれの座標頂点に取り、中央位置の直線を $y(t)$ とする。
制限 $w(-,y(t),-,-)$ の前後の座標 $i,j$ の成分は残る位置に値を持つ一次式で、
前の点との隣接から $t_i$ で、後の点との隣接から $t_j$ で零になる。
$i\ne j$ なら恒等零で、対角成分は $t_i$ で消える一次式と定数ベクトルの積である。
制限全体は零ではない（一文字縮約は非零）。非零の対角成分を前位置で取り出し
一般の $t$ を代入すれば、後位置から始まる関係空間に非零な純関係を得て矛盾する。
いまは $A,B$ による両偶座標が四頂点なのでまさにこの禁止配置である。
全ての場合が尽きたため奇側の反対対称性が従う。
\end{proof}
